\section*{A joint outer--child payment on the tensor reserve obstruction}

All logarithms are base two.  Let $u,v:\{-1,1\}^4\to\{-1,1\}$ be the
two Boolean slices from the independently accepted tensor-reserve no-go:
\[
u=(-1,1,-1,-1,-1,-1,1,-1,1,1,1,-1,1,-1,1,-1)
\]
and let $v$ be full parity.  Exact Walsh arithmetic gives
\[
 \langle u,v\rangle=\frac18,\qquad I(u)=\frac74,
 \qquad I(v)=4.
\]
For $k\ge1$, put $U_k=u^{\otimes k}$, $V_k=v^{\otimes k}$, and let
$d_{8k}$ be the point defect which is $+1$ at one point and $-1$
elsewhere.  Define the genuine Boolean function on $12k$ physical bits
\[
 F_k(x,z)=
 \begin{cases}
 U_k(z),&d_{8k}(x)=-1,\\
 V_k(z),&d_{8k}(x)=+1.
 \end{cases}
\]

\paragraph{Theorem.}
For every integer $k\ge1$,
\[
 H(F_k)\le \frac{1728}{595}I(F_k)<3I(F_k).
\]
Consequently the upper-set potential satisfies
\[
 U(F_k)\le \frac{1728}{595}I(F_k).
\]
If $\overline\Theta_k$ is the endpoint-centered biased-child cocycle for
the outer slices $(U_k,V_k)$ at the physical bias
$\mu_k=-1+2^{1-8k}$, then the complete terminal reserve has the joint
payment
\[
 \overline\Theta_k(\mu_k)
 \le \frac{1728}{595}I(F_k).
\]
The constant is independent of tensor depth, total dimension, bias,
cap-ratio growth, and coherent collision ratio.

\paragraph{Proof.}
Write $n=8k$, $r=2^{-n}$, $q=1-r$, and
\[
 \mu=-1+2r,\qquad \delta=1-\mu^2=4r(1-r).
\]
For the outer gate write $a=(U_k+V_k)/2$ and
$b=(V_k-U_k)/2$.  Its distinguished-input influence is
\[
 \beta_k=\|b\|_2^2=\frac{1-8^{-k}}2\le\frac12.
\]
Fourier coefficients of $F_k=a+d_n b$ split into two disjoint sectors:
\[
 \widehat F_k(\varnothing,T)=\widehat a(T)+\mu\widehat b(T),
 \qquad
 \widehat F_k(S,T)=\widehat d_n(S)\widehat b(T)\quad(S\ne\varnothing).
\]
The total squared mass of the nonempty-$x$ sector is
\[
 m=\delta\beta_k\le2r=2^{1-8k}.
\]
The empty sector has at most $2^{4k}$ atoms, and the nonempty sector has
fewer than $2^{12k}$ atoms.  Maximizing entropy within each sector gives
\[
 H(F_k)\le h_2(m)+4k(1-m)+12km
          =4k+h_2(m)+8km.
\]
Since $h_2(m)\le1\le k$ and
\[
8km\le16k\,2^{-8k}\le\frac{k}{16},
\]
we obtain
\[
 H(F_k)\le\frac{81}{16}k.                         \tag{1}
\]

Influence has the exact physical mixture identity
\[
 I(F_k)=qI(U_k)+rI(V_k)+\beta_k I(d_n).
\]
The tensor law for disjoint Boolean products gives
$I(U_k)=kI(u)=7k/4$.  As $q\ge255/256$,
\[
 I(F_k)\ge\frac{255}{256}\frac74 k
          =\frac{1785}{1024}k.                    \tag{2}
\]
Dividing (1) by (2) proves
\[
 \frac{H(F_k)}{I(F_k)}
 \le\frac{81/16}{1785/1024}
 =\frac{1728}{595}<3.
\]

For every spectral atom, its upper mass is at most one, so $U(f)\le H(f)$.
This proves the $U$ estimate.  Finally, the exact biased-child identity is
\[
 U(F_k)=\beta_kU(d_n)+\Theta_k(\mu_k),
\]
while $\Theta_k(-1)=U(U_k)$ and $\Theta_k(1)=U(V_k)$.  Hence
\[
 \overline\Theta_k(\mu_k)
 =U(F_k)-\beta_kU(d_n)-qU(U_k)-rU(V_k)
 \le U(F_k).
\]
Combining with the proved $U$ estimate finishes the joint payment.\qed

\paragraph{Scope.}
The child-only ratios from the preceding no-go still diverge: the theorem
does not repair either false child-only inequality.  It proves instead that
the exact same Boolean family has a uniform payment once the physical outer
budget is retained.  The proof uses the special fixed base slices and their
linear outer influence; it does not establish a universal joint ledger or
GU for arbitrary Boolean functions.
